% ============================================================================
% CHAPTER 16 APPENDIX: Domain-Specific HITL Metrics
% Measurement frameworks for audio, time-series, scientific discovery,
% and physical systems
% ============================================================================
% Source: /markdown/ch16/Ch16A_final.md
% ============================================================================

\chapter*{Technical Appendix 16A: Domain-Specific HITL Metrics}
\addcontentsline{toc}{chapter}{Technical Appendix 16A: Domain-Specific HITL Metrics}
\label{app:ch16}

% ============================================================================
\apxnumon
\section{Audio and Speech HITL Metrics}
% ============================================================================

\subsection{Word Error Rate with Demographic Slicing}

\term{Word Error Rate} (WER) is the standard transcription metric:

\begin{equation}
WER = \frac{S + D + I}{N}
\label{eq:wer}
\end{equation}

where $S$ is substitutions, $D$ is deletions, $I$ is insertions, and $N$ is the number of words in the reference transcript.

For \hitl audit purposes, WER must be computed by speaker demographics (not just in aggregate). A flagging threshold of 30\% worse in relative terms---that is, a slice WER more than $1.3\times$ the overall WER---is a practical starting point; the appropriate threshold depends on the application's sensitivity requirements.

\begin{technote}[WER Gap in Practice]
Koenecke et al. (2020) found that five major commercial \asr systems had an average WER of approximately 35\% for Black speakers versus 19\% for white speakers---roughly a 2$\times$ gap. This represents a systematic \hitl failure: annotation from these systems will have systematically different reliability across speaker demographics.
\end{technote}

\subsection{Diarization Error Rate}

\term{Diarization Error Rate} (DER) evaluates speaker attribution:

\begin{equation}
DER = \frac{\text{Missed Speech} + \text{False Alarm Speech}
            + \text{Speaker Confusion}}{\text{Total Reference Speech Time}}
\label{eq:der}
\end{equation}

For \hitl routing in legal and medical contexts, DER by confidence bin is critical: if high-confidence attributions are not more accurate than low-confidence attributions, the confidence score cannot be used to route cases to human review.

\subsection{Multimodal Crisis Escalation Logic}

The crisis detection escalation combines semantic and prosodic risk scores:

\begin{equation}
\text{Escalate} = \begin{cases}
1 & \text{if } r_{\text{combined}} > \theta_H \\
1 & \text{if } |r_{\text{text}} - r_{\text{prosody}}| > \delta \\
0 & \text{otherwise}
\end{cases}
\label{eq:crisis-escalation}
\end{equation}

where $r_{\text{combined}} = \alpha r_{\text{text}} + (1-\alpha) r_{\text{prosody}}$, $\theta_H$ is the high-risk threshold, and $\delta$ is the disagreement threshold. Setting $\delta$ appropriately ensures that exactly the kind of case this chapter opened with---text risk low, prosodic risk high---always triggers escalation.

% ============================================================================
\section{Time-Series HITL Metrics}
% ============================================================================

\subsection{Alarm Quality Metrics}

Four metrics characterize alarm system quality:

\paragraph{Alarm Precision.} Of all alarms generated, what fraction were genuine anomalies?
\begin{equation}
\text{Precision} = \frac{TP}{TP + FP}
\label{eq:alarm-precision}
\end{equation}

\paragraph{Alarm Burden.} Average alarms per operator per hour. EEMUA Publication 191 gives steady-state guidance: fewer than 6 alarms per hour (one per ten minutes) is very likely acceptable, 12 per hour (one per five minutes) is manageable, 30 per hour (one per two minutes) is likely over-demanding, and sustained rates above 60 per hour are very likely unacceptable \autocite{eemua191}.

\paragraph{Actionability Rate.} Fraction of reviewed alarms that resulted in a human intervention.

\paragraph{Critical Miss Rate.} Fraction of genuine anomalies that the system failed to alarm on.

\begin{table}[h]
\centering
\caption{Alarm Management Performance Thresholds (EEMUA 191)}
\label{tab:alarm-thresholds}
\begin{tabularx}{\textwidth}{@{}lYY@{}}
\toprule
\textbf{Alarm Burden} & \textbf{Category} & \textbf{HITL Implication} \\
\midrule
$< 6$/hr & Very likely acceptable & Target range for manageable oversight \\
$\sim 12$/hr & Manageable & Upper bound for sustained oversight \\
$\sim 30$/hr & Likely over-demanding & Alarm fatigue; redesign filtering \\
$> 60$/hr & Very likely unacceptable & Immediate intervention required \\
\bottomrule
\end{tabularx}
\end{table}

Rates well below one alarm per hour may instead indicate under-detection. Note also that these are steady-state averages: EEMUA 191 judges more than roughly 20 alarms in the first ten minutes after a major upset hard to cope with---flood-rate guidance, not a steady-state threshold.

\subsection{Temporal Uncertainty Quantification}

For time-series anomaly detection, uncertainty estimates must account for temporal context. Monte Carlo Dropout applied to a sequence model provides epistemic uncertainty estimates that can be augmented with changepoint proximity:

\begin{equation}
u_t^* = u_t \cdot (1 + \beta \cdot \text{CP}_t)
\label{eq:temporal-uncertainty}
\end{equation}

where $u_t$ is the base uncertainty estimate, $\text{CP}_t$ is a changepoint proximity score (0--1), and $\beta$ is a scaling factor. Uncertainty is higher near recent changepoints in the input signal.

% ============================================================================
\section{Scientific Discovery HITL Metrics}
% ============================================================================

\subsection{AlphaFold pLDDT Confidence Classification}

\begin{table}[h]
\centering
\caption{AlphaFold pLDDT Score Interpretation for Drug Design HITL}
\label{tab:plddt}
\begin{tabularx}{\textwidth}{@{}lYYY@{}}
\toprule
\textbf{pLDDT Range} & \textbf{Interpretation} & \textbf{Structural Reliability} & \textbf{Validation Required} \\
\midrule
$\geq 90$ & Very high & Comparable to experiment & Low \\
$70$--$90$ & High & Backbone reliable & Medium (side chains) \\
$50$--$70$ & Low & Uncertain & High \\
$< 50$ & Very low & Likely disordered & Critical \\
\bottomrule
\end{tabularx}
\end{table}

For drug design applications, the critical metric is the fraction of \textbf{binding site} residues in low or very-low confidence categories---not the overall structure mean. A protein with mean pLDDT of 85 may still have an unreliable binding site if the active site residues cluster in the 50--70 range.

\subsection{Bayesian Optimization Loop Metrics}

\begin{description}
    \item[Convergence rate] Number of iterations to reach 90\% of the best observed objective value
    \item[Human override rate] Fraction of model suggestions the researcher overrode
    \item[Human value-add] Average objective improvement of overridden experiments minus average objective improvement of model-suggested experiments. Positive values indicate the researcher is adding value beyond the model; negative values may indicate the researcher is underconfident in the model.
    \item[Data quality flag rate] Fraction of experimental outcomes flagged by the researcher as unreliable. High rates indicate an experimental process quality issue rather than a model quality issue.
\end{description}

% ============================================================================
\section{Robotics and Physical Systems HITL Metrics}
% ============================================================================

\subsection{Shared Autonomy Calibration}

A key metric for shared autonomy systems is the correlation between system uncertainty and human intervention rate:

\begin{equation}
\rho = \text{Corr}(\text{Uncertainty}_t, \mathbb{1}[\text{Human Intervened}_t])
\label{eq:shared-autonomy-corr}
\end{equation}

A high positive correlation ($\rho > 0.5$) indicates the system is correctly identifying when to ask for help. A low or negative correlation indicates that the system is asking for help on the wrong cases---a \textbf{Timing} failure in the Five Dimensions framework. Because most deployed systems trigger interventions by an uncertainty threshold, $\rho$ is high partly by construction, so it must be read jointly with the uncertainty--error correlation $\text{Corr}(\text{Uncertainty}_t, \mathbb{1}[\text{Model Error}_t])$: high $\rho$ with low uncertainty--error correlation means the system reliably asks for help on the \emph{wrong} cases.

\subsection{Prosthetic Adaptation Metrics}

\begin{description}
    \item[Adaptation trajectory] Accuracy slope per session (accuracy fraction per session). A positive slope indicates continuing improvement; slope $< 0.005$ suggests convergence.
    \item[Correction rate] Fraction of actuations where the user re-attempted due to incorrect movement. Above 15\% suggests recalibration is needed.
    \item[Per-movement accuracy] Accuracy by intended movement class. Identifies specific movements that require targeted retraining.
    \item[Latency] Time from muscle signal to actuation onset. Should remain stable or decrease as adaptation proceeds; increasing latency may indicate hardware issues.
\end{description}

% ============================================================================
\section{Universal Metrics for Non-Static HITL Systems}
% ============================================================================

\subsection{Distribution Shift Detection}

The Kolmogorov-Smirnov two-sample test detects whether the current data distribution has shifted from the training distribution:

\begin{equation}
D_{KS} = \sup_x |F_1(x) - F_2(x)|
\label{eq:ks-test}
\end{equation}

where $F_1$ and $F_2$ are the empirical CDFs of the training and recent data samples. Applied feature-by-feature, this identifies which input dimensions have drifted---pointing toward which aspect of the model's training distribution is no longer representative.

\subsection{HITL Health Dashboard Metrics}

For continuous deployment monitoring:

\begin{itemize}
    \item \textbf{Human-model agreement rate}: Fraction of cases where the human agrees with the model's recommendation. Below 70\% suggests model degradation or threshold misalignment.
    \item \textbf{Review rate}: Fraction of cases routed to human review. Above the target range suggests model degradation or threshold too conservative; below suggests overconfidence.
    \item \textbf{Reviewer decision consistency}: Cohen's $\kappa$ between the same reviewer's decisions on matched cases at different time points. Below 0.6 suggests reviewer fatigue or changing standards.
    \item \textbf{Outcome tracking}: Where outcomes are observable, track accuracy of decisions including human review, not just pre-review model predictions.
\end{itemize}
\apxnumoff
